% ---------------------------------------------------------------------
% 第 5 章 几何一致控制律与认证 H∞ 性能
% ---------------------------------------------------------------------
\section{几何一致控制律与认证 H∞ 性能}

\subsection{扰动项：显式解算、乘性项分离与假设集}

真实关节空间动力学（含摩擦 $\boldsymbol f$、执行器实现误差 $\delta\boldsymbol\tau$ 与环境接触 $\boldsymbol\tau_{\mathrm{ext}}$）记为 $M(\boldsymbol q)\ddot{\boldsymbol q}+C(\boldsymbol q,\dot{\boldsymbol q})\dot{\boldsymbol q}+\boldsymbol g(\boldsymbol q)+\boldsymbol f=\boldsymbol\tau+\delta\boldsymbol\tau+\boldsymbol\tau_{\mathrm{ext}}$；内环以标称模型执行计算力矩 $\boldsymbol\tau=\hat M\ddot{\boldsymbol q}_{\mathrm{ref}}+\hat C\dot{\boldsymbol q}+\hat g+\hat{\boldsymbol f}$。失配量取“标称减真实”：$\Delta M\triangleq\hat M-M$，$\Delta C,\Delta\boldsymbol g,\Delta\boldsymbol f$ 同理。代入并左乘 $M^{-1}$ 得 $\ddot{\boldsymbol q}=\ddot{\boldsymbol q}_{\mathrm{ref}}+\boldsymbol w_{\mathrm{dyn}}$，其中
\begin{equation}
\boldsymbol w_{\mathrm{dyn}}=M^{-1}\bigl(\Delta M\,\ddot{\boldsymbol q}_{\mathrm{ref}}+\Delta C\dot{\boldsymbol q}+\Delta\boldsymbol g+\Delta\boldsymbol f+\delta\boldsymbol\tau+\boldsymbol\tau_{\mathrm{ext}}\bigr).
\tag{5.1}\label{eq:5.1}
\end{equation}

\textbf{反馈乘性项与余项分离}。记
\[
u_{\rm ff}\triangleq\mathrm{vec}_6\!\left(\mathrm{Ad}_{\tilde x}\dot{\boldsymbol\xi}_d+\mathrm{ad}_{\tilde{\boldsymbol\xi}}\mathrm{Ad}_{\tilde x}\boldsymbol\xi_d\right),\qquad
u_{\rm fb}\triangleq-K_de_\xi-A^{\top}(\tilde x)K_pe_z .
\]
将 $\ddot{\boldsymbol q}_{\mathrm{ref}}=J^{+}(u_{\rm ff}+u_{\rm fb}-\dot J\dot{\boldsymbol q})$ 代入 \eqref{eq:5.1}，利用 $JJ^{+}=I_6$ 得扰动分解
\begin{equation}
d=\Theta(\boldsymbol q)\,u_{\mathrm{fb}}+d_{\mathrm{ex}},\qquad
\Theta\triangleq J M^{-1}\Delta M\,J^{+}\in\mathbb R^{6\times6},\;
\tag{5.1d}\label{eq:5.1d}
\end{equation}
其中 $\Theta$ 为反馈乘性分量，$d_{\mathrm{ex}}$ 为含状态相关项的余项（附录 C.1）。定理 3(c) 假设总扰动 $d\in L_2$，3(d) 则假设余项有界。

\textbf{假设集}（各结论按需引用）：
\begin{itemize}[leftmargin=2em]
  \item \textbf{(A1) 雅可比正则}：在工作集 $\mathcal Q$ 上 $\sigma_{\min}(J)\ge\underline\sigma>0$，理论取 Moore--Penrose 伪逆 $J^{+}$，故 $JJ^{+}=I_6$。阻尼伪逆产生的 $JJ^{\#}-I_6$ 残差不在此精确实现模型内。
  \item \textbf{(A2) 惯量与失配有界}：$M(\boldsymbol q)\succeq \underline m I_n$（$\underline m>0$），且 $\Delta M,\Delta C,\Delta\boldsymbol g,\Delta\boldsymbol f$ 在 $\mathcal Q\times\{\|\dot{\boldsymbol q}\|\le\bar v\}$ 上有界。
  \item \textbf{(A3) 乘性小增益条件}：
  \begin{equation}
  \alpha\triangleq\sup_{\boldsymbol q\in\mathcal Q}\bigl\|\Theta(\boldsymbol q)\bigr\|_2
  \;<\;\frac{\lambda_{\min}(K_d)}{\lambda_{\max}(K_d)}=\frac1{\mathrm{cond}_2(K_d)}\ \le 1 .
  \tag{5.1f}\label{eq:5.1f}
  \end{equation}
  各向同性阻尼时退化为 $\alpha<1$。
  \item \textbf{(A4) 轨迹与扰动正则性}：$\hat{\underline x}_d(t)$ 为 $C^2$ 且 $\|\boldsymbol\xi_d\|,\|\dot{\boldsymbol\xi}_d\|$ 有界。3(c) 要求总扰动 $d\in L_2$；3(d) 要求 $d_{\mathrm{ex}}\in L_\infty$，二者为不同情形。
\end{itemize}

\subsection{控制律}

取 $K_p$ 对称正定（典型取块对角 $K_p=\mathrm{diag}(p_O I_3,p_T I_3)$；取 $K_p=k_pI_6$ 即恢复单标量增益情形）、$K_d$ 对称正定（若需定理 3(c-2) 的旋转/平移逐分量 H∞ 界，则取块对角 $K_d=\mathrm{diag}(K_\omega,K_v)$，$K_\omega,K_v\in\mathbb R^{3\times3}$ 对称正定），定义参考加速度（几何一致计算力矩律）：
\begin{equation}
\ddot{\boldsymbol q}_{\mathrm{ref}}
=J^{+}\Bigl(\underbrace{\mathrm{vec}_6\bigl(\mathrm{Ad}_{\tilde x}\dot{\boldsymbol\xi}_d+\mathrm{ad}_{\tilde{\boldsymbol\xi}}\mathrm{Ad}_{\tilde x}\boldsymbol\xi_d\bigr)}_{\text{前馈：搬运的期望加速度 + 输运修正}}
-K_d\,e_\xi-A^{\top}(\tilde x)\,K_p\,e_z-\dot J\dot{\boldsymbol q}\Bigr).
\tag{5.2}\label{eq:5.2}
\end{equation}
其中 $A^\top$-整形位姿反馈使 Lyapunov 导数的交叉项精确相消（定理 3 证明），前馈与 $\dot J\dot{\boldsymbol q}$ 分别取自期望链的 $\sigma^2$ 项与 \eqref{eq:3.5}；全部反馈量取自定理 1 的 HDQ 误差元素，不需要任何加速度误差的测量或估计。

\subsection{主定理}

以下分析理想连续误差模型。机械臂实现另假设解在所考察时段内存在并留在 (A1)、(A2) 的正则工作集，期望轨迹满足 (A4)；无穷时域结论要求该前提全时成立，不能由任务误差有界推出。四元数采用连续提升，仅在初始时选符号。统一存储函数为
\begin{equation}
V(e_z,e_\xi)=\tfrac12\|e_\xi\|^2+\tfrac12\,e_z^{\top}K_pe_z\;\ge0 ,
\tag{5.4a}\label{eq:5.4a}
\end{equation}
$K_p\succ0$ 对称。无扰收敛限于正四元数分支的水平集；耗散不等式不要求 $A$ 可逆。

% ---- 定理 3(a)：手工编号 ----
\setcounter{theorem}{2}%
\renewcommand{\thetheorem}{3(a)}%
\begin{theorem}[闭环误差动态]\label{thm:3a}
误差坐标 $(e_z,e_\xi)$ 满足二阶系统
\begin{equation}
\dot e_z=A(\tilde x)\,e_\xi,\qquad
\dot e_\xi=-K_d e_\xi-A^{\top}(\tilde x)\,K_p\,e_z+d(t),
\tag{5.5}\label{eq:5.5}
\end{equation}
其中 $d=J\boldsymbol w_{\mathrm{dyn}}$，并按 \eqref{eq:5.1d} 分解为 $d=\Theta u_{\mathrm{fb}}+d_{\mathrm{ex}}$。
\end{theorem}

\begin{proof}
对 \eqref{eq:4.4} 逐项求导，用伴随求导法则 $\frac{d}{dt}\bigl(\mathrm{Ad}_{\tilde x}\boldsymbol a\bigr)=\mathrm{Ad}_{\tilde x}\dot{\boldsymbol a}+\mathrm{ad}_{\tilde{\boldsymbol\xi}}(\mathrm{Ad}_{\tilde x}\boldsymbol a)$（由 $\dot{\tilde x}=\tfrac12\tilde{\boldsymbol\xi}\tilde x$ 与其共轭直接展开）得 twist 误差动态
\begin{equation}
\dot{\tilde{\boldsymbol\xi}}=\dot{\boldsymbol\xi}-\mathrm{Ad}_{\tilde x}\dot{\boldsymbol\xi}_d-\mathrm{ad}_{\tilde{\boldsymbol\xi}}\bigl(\mathrm{Ad}_{\tilde x}\boldsymbol\xi_d\bigr)
\tag{5.5a}\label{eq:5.5a}
\end{equation}
将 \eqref{eq:5.2} 与 $\ddot{\boldsymbol q}=\ddot{\boldsymbol q}_{\mathrm{ref}}+\boldsymbol w_{\mathrm{dyn}}$ 代入 \eqref{eq:3.5}，再代入上式，前馈相消，余下 $d=J\boldsymbol w_{\mathrm{dyn}}$ 即得 \eqref{eq:5.5}。
\end{proof}

% ---- 定理 3(b) ----
\renewcommand{\thetheorem}{3(b)}%
\begin{theorem}[无扰：水平集不变性、渐近收敛与局部指数稳定]\label{thm:3b}
设 $d\equiv0$，$K_p=\mathrm{diag}(K_{p,O},K_{p,T})$ 对称正定（旋转/平移块）。取
\begin{equation}
0<c<c^{*}\triangleq\tfrac12\lambda_{\min}(K_{p,O}),
\qquad
\Omega_c\triangleq\{(e_z,e_\xi):V\le c\}.
\tag{5.5b}\label{eq:5.5b}
\end{equation}
则：\textbf{(i)} $\dot V=-e_\xi^\top K_de_\xi\le0$，故 $\Omega_c$ 紧且正向不变；\textbf{(ii)} 初值在 $\Omega_c$ 内且 $\tilde\eta(0)>0$ 时，全程有
\begin{equation}
\tilde\eta(t)\ \ge\ \eta_0\triangleq\sqrt{1-\frac{2c}{\lambda_{\min}(K_{p,O})}}>0 ,
\tag{5.5c}\label{eq:5.5c}
\end{equation}
从而 $A(\tilde x)$ 在 $\Omega_c$ 上一致可逆（$\det A=-\tfrac18\tilde\eta\le-\tfrac18\eta_0<0$）；\textbf{(iii)} $\Omega_c$ 内一切轨迹满足 $(e_z,e_\xi)\to(0,0)$；\textbf{(iv)} $(0,0)$ 是\textbf{局部指数稳定}的。
\end{theorem}

\begin{proof}
关键步骤是 $A^\top$-整形反馈使交叉项 $e_z^\top K_pAe_\xi$ 精确相消，得 $\dot V=-e_\xi^\top K_de_\xi$；LaSalle 定理与 $A$ 在 $\Omega_c$ 上的一致可逆性（附录 C.2）给出渐近收敛。详见附录 C.4。
\end{proof}

\begin{remark}
初始 $\tilde\eta<0$ 时可将 HDQ 误差整体翻号，不改变 $\tilde{\boldsymbol\xi}$；上述水平集保证后续无需切换。$\blacksquare$
\end{remark}

% ---- 定理 3(c) ----
\renewcommand{\thetheorem}{3(c)}%
\begin{theorem}[$L_2$ 扰动：H∞ 二次型/Schur 补判据与旋转/平移分量拆分]\label{thm:3c}
设 $d=d_{L_2}\in L_2$。

\textbf{证书参数}：$\kappa,\gamma_a>0$ 不进入控制律，认证 $d\to e_\xi$ 的增益 $\gamma^{*}=\gamma_a\sqrt\kappa$（量纲 s）；零初值下 $\|e_\xi\|_{L_2}\le\gamma^{*}\|d\|_{L_2}$。证书族内的阻尼反演见 3(c$'$)。

\textbf{(c-1) 合并判据}（$K_d$ 任意对称正定）：对给定存储函数与供给率，点态耗散不等式 $\dot V\le-\tfrac1{2\kappa}\|e_\xi\|^2+\tfrac{\gamma_a^2}2\|d\|^2$ 对任意 $(e_\xi,d)$ 成立的充要条件为
\begin{equation}
M\triangleq\begin{bmatrix}K_d-\tfrac1{2\kappa}I_6 & -\tfrac12 I_6\\[2pt] -\tfrac12 I_6 & \tfrac{\gamma_a^2}2 I_6\end{bmatrix}\succeq0
\;\overset{\text{Schur}}{\Longleftrightarrow}\;
K_d\succeq\tfrac12\bigl(\kappa^{-1}+\gamma_a^{-2}\bigr)I_6 ,
\tag{5.6a}\label{eq:5.6a}
\end{equation}
此时
\begin{equation}
\int_0^\infty\kappa^{-1}\|e_\xi\|^2dt\le\gamma_a^2\int_0^\infty\|d_{L_2}\|^2dt+2V(0),
\tag{5.6}\label{eq:5.6}
\end{equation}
即加速度层扰动到 twist 误差能量的 $L_2$ 增益 $\le\gamma_a\sqrt\kappa$（零初值时退化为纯增益界；此界只约束 $e_\xi$）。

\textbf{(c-2) 分量拆分判据}（$K_d=\mathrm{diag}(K_\omega,K_v)$ 块对角，且 $K_p=\mathrm{diag}(K_{p,O},\,k_{p,T}I_3)$，平移刚度块须各向同性，必要性见附录 C.3）：记 $e_\xi=[\tilde\omega;\tilde v]$、$d=[d_\omega;d_v]$，$V_\omega\triangleq\tfrac12\|\tilde\omega\|^2+\tfrac12\mathcal O^\top K_{p,O}\mathcal O$、$V_v\triangleq\tfrac12\|\tilde v\|^2+\tfrac{k_{p,T}}2\|\mathcal T\|^2$（$V=V_\omega+V_v$）。若
\begin{equation}
K_\omega\succeq\tfrac12\bigl(\kappa_\omega^{-1}+\gamma_\omega^{-2}\bigr)I_3,
\qquad
K_v\succeq\tfrac12\bigl(\kappa_v^{-1}+\gamma_v^{-2}\bigr)I_3,
\tag{5.6b}\label{eq:5.6b}
\end{equation}
则旋转/平移两分量\textbf{各自独立}满足
\begin{equation}
\int_0^\infty\kappa_\omega^{-1}\|\tilde\omega\|^2dt\le\gamma_\omega^2\int_0^\infty\|d_\omega\|^2dt+2V_\omega(0),
\quad
\int_0^\infty\kappa_v^{-1}\|\tilde v\|^2dt\le\gamma_v^2\int_0^\infty\|d_v\|^2dt+2V_v(0),
\tag{5.6$'$}\label{eq:5.6p}
\end{equation}
两组证书参数可独立指定。
\end{theorem}

\begin{proof}
含扰时交叉项仍精确相消，得 $\dot V=-e_\xi^\top K_de_\xi+e_\xi^\top d$；将性能目标写为 $[e_\xi;d]^\top M[e_\xi;d]\ge0$，取 Schur 补即得充要判据 \eqref{eq:5.6a}；分量解耦依赖两处代数恒零（$\mathcal T\times\mathcal T=0$、$\mathcal T\cdot(\mathcal T\times\tilde\omega)=0$）。详见附录 C.5。
\end{proof}

% ---- 定理 3(c')：最优证书与 γ 反演 ----
\renewcommand{\thetheorem}{3(c$'$)}%
\begin{theorem}[最优证书与目标增益反演：$\gamma^{*}$ 调参]\label{thm:3c2}
记认证 $L_2$ 增益 $\gamma^{*}\triangleq\gamma_a\sqrt{\kappa}$，$\lambda\triangleq\lambda_{\min}(K_d)>0$。则：
\textbf{(i)}（下界）证书族 \eqref{eq:5.6a} 的任一可行点都满足 $\gamma^{*}\ge1/\lambda$；
\textbf{(ii)}（可达）$(\kappa,\gamma_a)=(\lambda^{-1},\lambda^{-1/2})$ 可行且 $\gamma^{*}=1/\lambda$。
故证书族可认证的最优（最小）增益为
\begin{equation}
\gamma^{*}_{\mathrm{opt}}\;=\;\min_{(\kappa,\gamma_a)\ \mathrm{feasible}}\gamma_a\sqrt{\kappa}\;=\;\frac{1}{\lambda_{\min}(K_d)} ,
\tag{5.6c}\label{eq:5.6c}
\end{equation}
最优点为 $\kappa=\gamma_a^2=1/\lambda_{\min}(K_d)$。\textbf{在该证书族内}，认证目标 $\gamma^{*}$ 当且仅当 $K_d\succeq(1/\gamma^{*})I_6$，可行点为 $(\kappa,\gamma_a)=(\gamma^{*},\sqrt{\gamma^{*}})$。式 \eqref{eq:5.6b} 同理给出最优分量界 $\gamma^{*}_\omega=1/\lambda_{\min}(K_\omega)$、$\gamma^{*}_v=1/\lambda_{\min}(K_v)$。
\end{theorem}

\begin{proof}
(i)~由 AM--GM 不等式（各项为正），$\kappa^{-1}+\gamma_a^{-2}\ \ge\ 2\sqrt{\kappa^{-1}\gamma_a^{-2}}\ =\ \dfrac{2}{\gamma_a\sqrt{\kappa}}\ =\ \dfrac{2}{\gamma^{*}}$；代入判据 \eqref{eq:5.6a} 得 $\lambda\ \ge\ \tfrac12\bigl(\kappa^{-1}+\gamma_a^{-2}\bigr)\ \ge\ 1/\gamma^{*}$。
(ii)~在 $\kappa=\gamma_a^2=\lambda^{-1}$ 处 $\kappa^{-1}+\gamma_a^{-2}=2\lambda$，\eqref{eq:5.6a} 取等成立，证书增益 $\gamma_a\sqrt{\kappa}=1/\lambda$；且 $\kappa=\lambda^{-1}>1/(2\lambda)$，可行域非空。
反演规则的必要性即 (i)；充分性：若 $\lambda\ge1/\gamma^{*}$，则点 $(\kappa,\gamma_a)=(\gamma^{*},\sqrt{\gamma^{*}})$ 满足 $\kappa^{-1}+\gamma_a^{-2}=2/\gamma^{*}\le2\lambda$，可行，且其证书增益恰为 $\gamma_a\sqrt{\kappa}=\gamma^{*}$。
\end{proof}

\begin{remark}
当 $K_p=\mathrm{diag}(p_OI_3,p_TI_3)$、$K_d=\mathrm{diag}(K_\omega,K_v)$ 时，认证值等于线性化 $H_\infty$ 范数；一般 SPD 增益下仅为上界（附录 C.7）。$\blacksquare$
\end{remark}

% ---- 定理 3(d) ----
\renewcommand{\thetheorem}{3(d)}%
\begin{theorem}[$L_\infty$ 扰动：乘性分量分离与 twist 误差的均方极限界]\label{thm:3d}
设 $K_p=\mathrm{diag}(K_{p,O},K_{p,T})\succ0$，(A3) 成立，扰动按 \eqref{eq:5.1d} 分解，$D_{\mathrm{ex}}\triangleq\|d_{\mathrm{ex}}\|_{L_\infty}<\infty$。记\textbf{有效阻尼}
\begin{equation}
\lambda_{\mathrm{eff}}\triangleq\lambda_{\min}(K_d)-\alpha\,\lambda_{\max}(K_d)\;>\;0
\tag{5.7a}\label{eq:5.7a}
\end{equation}
设轨迹在 $[0,T_\Omega]$ 内留在 $\Omega_c$（\eqref{eq:5.5b}），记\textbf{等效扰动幅值}
\begin{equation}
D\triangleq D_{\mathrm{ex}}+\alpha\,\lambda_{\max}(K_p)\Bigl(1+\sqrt{\tfrac{2c}{\lambda_{\min}(K_{p,T})}}\Bigr)\sqrt{\tfrac{2c}{\lambda_{\min}(K_p)}} .
\tag{5.7b}\label{eq:5.7b}
\end{equation}
则 twist 误差满足有限时段均方界
\begin{equation}
\frac1T\int_0^T\|e_\xi(t)\|^2\,dt\;\le\;\frac{D^2}{\lambda_{\mathrm{eff}}^2}+\frac{2V(0)}{\lambda_{\mathrm{eff}}\,T},
\qquad 0<T\le T_\Omega,
\tag{5.7c}\label{eq:5.7c}
\end{equation}
若该水平集前提对所有 $t\ge0$ 成立，则进一步有
\begin{equation}
\limsup_{T\to\infty}\ \mathrm{RMS}_{[0,T]}(e_\xi)
\;\triangleq\;\limsup_{T\to\infty}\Bigl(\frac1T\int_0^T\|e_\xi\|^2dt\Bigr)^{1/2}
\;\le\;\frac{D}{\lambda_{\mathrm{eff}}}
\;\xrightarrow[\ \alpha\to0\ ]{}\;\frac{\|d_{\mathrm{ex}}\|_{L_\infty}}{\lambda_{\min}(K_d)} .
\tag{5.7}\label{eq:5.7}
\end{equation}
这是以留域为前提的 twist 均方界，不证明受扰水平集不变性、全状态 ISS 或位姿稳态标度律。
\end{theorem}

\begin{proof}
将乘性扰动 $\Theta u_{\mathrm{fb}}$ 拆为阻尼折减项与等效扰动项，分别用 Young 不等式回收，积分即得 RMS 界 \eqref{eq:5.7}。详见附录 C.6。
\end{proof}
